% !TEX root = ../../main_performance.tex
% ---------------------------------------------------------------------
%  Rekonfigurierbare Hardware
% ---------------------------------------------------------------------
\section{Rekonfigurierbare Hardware}
\label{sec:reconfig}

\SW{Spektrum} fein\,\ra\,grob, schnell\,\ra\,flexibel:
\begin{ztab}{@{}lXl@{}}
Ebene & Baustein & Eigenschaft \\ \midrule
Gatter & \T{FPGA} & sehr feingranular \\
Sequencer & progr.\ FSM & kaum progr., sehr schnell \\
 & Controller & winziger Befehlssatz (Load, Store, Add,
   Jump, Compare), flexibel, langsam \\
 & \T{PIO} & flexibel \emph{und} schnell \\
 & kleine CPU $+$ RAM & sehr flexibel, sehr langsam \\
Operation & \T{CGRA} & routbare ALU/FPU-Blöcke, Datenpfad, VLIW \\
\end{ztab}

\subsection{FPGA}
\begin{itemize}
\item \T{LUT} (Look Up Table) realisiert \emph{jede} Logikfunktion ihrer
      Eingänge; im Beispiel 3, real 4--8
\item Konfigurationsspeicher $=$ SRAM-Zellen in Schieberegister-Anordnung
      (Wiederverwendung der Scan-Test-Struktur)
\item Switch-Matrix-Routing mit Sechs-Transistor-Kreuzungen
\item konfiguriert in VHDL
\end{itemize}
\begin{formel}[Die Zahl, die man sich merkt]
\T{Routing-Delay $>$70\,\% des Gesamt-Delays.} Nicht die Logik ist
langsam, sondern die Verdrahtung.
\end{formel}

\subsection{CGRA}
Array aus \T{PE} (Processing Elements) mit \T{FU} (spezialisierte
ALU/FPU: AND, SHL, ADD, MUL, MAC), typisch extremes \T{VLIW}.

\begin{bsp}[4-Tap-FIR auf $4\times4$-CGRA]
\label{sec:cgra}
$y[n]=h_0x[n]+h_1x[n{-}1]+h_2x[n{-}2]+h_3x[n{-}3]$,
Nearest-Neighbor-Mesh, 1-Takt-ADD/MUL.
\[m_k=h_k\,x[n{-}k],\quad a_0=m_0{+}m_1,\ a_1=m_2{+}m_3,\ y=a_0{+}a_1\]
\T{7 von 16 PE} benutzt $\;\;$ \T{1 Ergebnis pro Takt} (Pipeline)
$\;\;$ Latenz $=1_{\text{MUL}}+2_{\text{ADD}}=\T{3 Takte}$.
\end{bsp}

\SW{Xilinx AI Engines} Kachel aus Vektorprozessor (SIMD für parallele
Arithmetik) $+$ Skalarprozessor (Kontrollfluss, Adressen) $+$ lokalem
deterministischem Speicher $+$ DMA-Engines, verbunden über streaming-
und memory-mapped Interconnect.

\subsection{RP2350 PIO}
3 identische PIO-Blöcke, je \T{4 State Machines}. Realisiert in Software
Schnittstellen wie I2C, 3-pin I2S, SDIO, SPI/DSPI/QSPI, UART, DPI,
8080/6800-Parallelbus, VGA (Widerstands-DAC).
\begin{ztab}{@{}lX@{}}
Register je SM & $2\times32$\,Bit Shift (\T{OSR}, \T{ISR}),
  $2\times32$\,Bit Scratch (\T{X}, \T{Y}) \\
FIFO & $4\times32$\,Bit je TX und RX \\
Instruktionsspeicher & 32 Worte, gemeinsam für die 4 SM \\
Takt & fraktionaler Teiler (16 Integer-, 8 Nachkommabits) \\
Rest & flexibles GPIO-Mapping, DMA-Interface, IRQ-Flags \\
\end{ztab}

\begin{formel}[Nur 9 Befehle]
\R{JMP} \R{WAIT} \R{IN} \R{OUT} \R{PUSH} \R{PULL} \R{MOV} \R{IRQ}
\R{SET}. \R{NOP} ist eine Pseudo-Instruktion ($=$ \R{mov y,y}).\\
Format: \R{<instr> (side <wert>) ([<delay>])}, Delay 0--31 Takte
(weniger, wenn \R{.side\_set} benutzt wird).\\
JMP-Bedingungen: immer, \R{!X}, \R{X--}, \R{!Y}, \R{Y--}, \R{X!=Y},
\R{PIN}, \R{!OSRE}.
\end{formel}

\SW{Datenfluss} \R{PULL} holt ein Wort aus der TX-FIFO in den \T{OSR},
\R{OUT} schiebt 1--32\,Bit daraus zum Ziel; \R{IN} schiebt in den
\T{ISR}, \R{PUSH} schreibt ihn in die RX-FIFO.
\T{Autopull/Autopush} füllt bzw.\ leert automatisch
\ra die \R{pull}/\R{push}-Instruktion entfällt.

\SW{Zwei Tricks, die Takte sparen}
\begin{description}
\item[Side-Set] schaltet \emph{parallel} zum Befehl Pins um
  \ra Daten \emph{und} Takt ohne Zusatzbefehl (SPI).
\item[Program Wrapping] \R{.wrap\_target}/\R{.wrap} ersetzt den \R{jmp}
  durch einen Rücksprung, der \emph{null} Takte kostet.
\end{description}

\begin{lstlisting}
.program squarewave_wrap
  set pindirs, 1     ; Pin auf Ausgang
.wrap_target         ; 0-Takt-Ruecksprungziel
  set pins, 1 [1]    ; high, 1 Takt warten
  set pins, 0 [1]    ; low,  1 Takt warten
.wrap                ; 0-Takt-Sprung statt jmp
\end{lstlisting}

\SW{Direktiven} \R{.program}, \R{.define}, \R{.side\_set <n> (opt)},
\R{.in}/\R{.out <n> (left|right) (auto) (<schwelle>)},
\R{.fifo}, \R{.clock\_div}, \R{.mov\_status},
\R{.wrap\_target}/\R{.wrap}, \R{.origin}, \R{.pio\_version}
(1 $=$ RP2350, 0 $=$ RP2040), \R{.word}.

\subsection{PIO-Programmiermodell}
\label{rez:pio}
\begin{ztab}{@{}lX@{}}
\multicolumn{2}{@{}l}{\textbf{Die 9 Instruktionen}}\\
\R{JMP} & \R{PC} setzen, bedingt auf \R{X}, \R{Y} oder Pin \\
\R{WAIT} & auf Pin, IRQ oder Clock warten -- blockierend \\
\R{IN} / \R{OUT} & 1--32 Bit zwischen \R{ISR}/\R{OSR} und Pins, \R{X},
  \R{Y} schieben \\
\R{PUSH} / \R{PULL} & 32-Bit-Wort \R{ISR}\,\ra\,RX-FIFO bzw.\
  TX-FIFO\,\ra\,\R{OSR} \\
\R{MOV} & zwischen Registern kopieren, optional invertieren/spiegeln \\
\R{IRQ} & Flag setzen/löschen -- die Synchronisation zwischen SMs und
  zur CPU \\
\R{SET} & Konstante in Pins, Pin-Richtung, \R{X} oder \R{Y} \\
\end{ztab}
\SW{Register pro State Machine} \R{PC}, \R{OSR} (Output Shift Register),
\R{ISR} (Input Shift Register), die Scratch-Register \R{X} und \R{Y},
dazu TX- und RX-FIFO mit je 4 Worten. Vier SMs teilen sich einen
\emph{gemeinsamen} Instruktionsspeicher von 32 Worten.

\begin{bsp}[Rechteck 50/50, Periode 4 Takte]
\begin{lstlisting}
.program squarewave
  set pindirs, 1    ; Pin als Ausgang
again:
  set pins, 1 [1]   ; high, dann 1 Takt warten
  set pins, 0       ; low
  jmp again         ; PC = again
\end{lstlisting}
Jede Instruktion braucht 1 Takt; \R{[1]} hängt einen Wartetakt an
\ra $1+1+1+1 = 4$ Takte, davon 2 high und 2 low. Takt teilen mit
\R{sm[0].clkdiv = (uint32\_t)(2.5f * (1 << 16))} (16.16-Festkomma).
\end{bsp}

\begin{rezept}[Ein PIO-Programm in Betrieb nehmen]
\begin{enumerate}
\item Programm in den Instruktionsspeicher schreiben:
      \R{pio->instr\_mem[i] = ...\_program\_instructions[i]}
\item \R{sm[n].clkdiv} setzen (Teiler als 16.16-Festkomma)
\item \R{sm[n].pinctrl} setzen: \R{SET\_COUNT} und \R{SET\_BASE} sagen,
      \emph{wie viele} Pins ab \emph{welchem} Pin die \R{SET}-Befehle
      treffen
\item GPIO auf die PIO-Funktion umschalten:
      \R{gpio\_set\_function(0, pio\_get\_funcsel(pio))}
\item SM freigeben: \R{hw\_set\_bits(\&pio->ctrl, 1 << (SM\_ENABLE + n))}
\end{enumerate}
\end{rezept}

\SW{\R{.wrap}} \R{.wrap\_target} und \R{.wrap} lassen die SM in Hardware
zurückspringen -- spart den \R{jmp} und damit einen Takt pro
Schleifendurchlauf. \R{side-set} legt gleichzeitig zur Instruktion Pins
um, ebenfalls ohne eigenen Takt; so entstehen exakte Taktsignale.

\begin{falle}[Warum PIO und nicht Bit-Banging]
Bit-Banging auf dem Core ist jitterbehaftet (Interrupts!) und blockiert
die CPU. Die SM läuft taktgenau und unabhängig, die CPU füttert nur die
FIFOs -- oder lässt das den DMA machen. Grenze: 32 Instruktionen,
keine Arithmetik ausser Dekrementieren, und zwei Scratch-Register.
\end{falle}

\subsection{Die drei Körnungen im Vergleich}
\label{rez:koernung}
\begin{ztab}{@{}lXX@{}}
 & Baustein & Eigenschaft \\ \midrule
\T{Fine Grain} & FPGA: LUT auf \emph{Bit}-Ebene & maximal flexibel,
  Routing dominiert \\
\T{Coarse Grain} & CGRA: ALU/FPU auf \emph{Wort}-Ebene & weniger
  flexibel, dafür schneller und kleiner \\
\T{Prozessor} & fester Befehlssatz & maximal flexibel in Software,
  langsam pro Operation \\
\end{ztab}
\SW{Dazwischen} \T{PIO} (Programmable I/O): flexibel \emph{und} schnell,
aber auf I/O-Muster beschränkt. \T{Kleine CPUs $+$ RAM} im Chip: sehr
flexibel, aber sehr langsam.

\begin{formel}[Warum das Routing das Problem ist]
Im FPGA entfallen \T{über 70\,\% der Gesamtverzögerung} auf das
\emph{Routing}, nicht auf die Logik. Die Switch-Matrix besteht aus
Six-Transistor-Kreuzungen; jedes Durchqueren kostet. Deshalb:
\begin{itemize}
\item die Taktfrequenz eines FPGA-Designs steht und fällt mit der
      Platzierung, nicht mit der Logiktiefe
\item ein CGRA ist schneller, weil seine Blöcke grösser sind und deshalb
      \emph{weniger Kreuzungen} pro Operation nötig sind
\end{itemize}
\end{formel}
\SW{LUT} Eine Look-Up-Table mit $n$ Eingängen kann \emph{jede} logische
Funktion dieser $n$ Eingänge -- sie ist einfach eine Wahrheitstabelle im
SRAM. Lehrbuchbild: 3 Eingänge, reale FPGAs 4--8. Der
Konfigurationsspeicher sind SRAM-Zellen in Schieberegister-Anordnung
(wiederverwendete Scan-Test-Struktur) -- deshalb ist ein FPGA nach dem
Einschalten leer und muss geladen werden.

\begin{falle}[Was „rekonfigurierbar“ nicht heisst]
Ein FPGA rechnet nicht schneller, weil er „Hardware“ ist, sondern weil
er \emph{parallel} und \emph{ohne Befehlsholen} arbeitet. Ein einzelner
LUT-Pfad ist \emph{langsamer} als eine CPU-ALU. Der Gewinn kommt aus der
Anzahl, nicht aus der Geschwindigkeit -- dieselbe Logik wie bei SIMD und
Multicore.
\end{falle}
